% MAS1004 Assignment 1: long report.
%
% You fill this in with your agent, in your repository. AGENTS.md tells the
% agent how. Red text in the PDF is what still has to be replaced.
%
% - Pictures go in figures/. Copy them from results/, which is not in git.
% - Outputs go in outputs/. run.py saves what it prints in
%   results/<tag>_output.txt, and PLAN.md saves the other outputs in results/
%   with `| tee results/...`. Copy them here. \outputfile{outputs/check.txt}
%   prints a file exactly as it is. Every number in this report has to come
%   from such an output.
%
% To make the PDF: compress this folder into a zip, and on Overleaf choose
% New Project, Upload Project. Set Menu, Compiler to XeLaTeX, and Recompile.

\documentclass[11pt]{article}
\usepackage{mas1004}

\reporttitle{Assignment 1: Long report}
\studentname{Hai Anh}
\studentid{ADD YOUR STUDENT ID}

\begin{document}
\makeheader

% ---------------------------------------------------------------------------
\section*{Problem 1. Categories and images}

\subsection*{Why these categories}
I chose cartoon animals, plush animals, and real animals because they let me compare three visual styles around a subject I encounter often in media and online images. The categories are related enough to make the task meaningful, but their appearance can overlap: a realistic plush toy may resemble a real animal, and a cartoon can share the shape and colours of either.

\subsection*{What I expected the model to use}
\fillin{What you thought, before training, the model would use to tell them
apart: shape, colour, texture, background, or something else.}

\subsection*{Images per category}
I downloaded 100 images per category (300 total). After removing 11 images, the cleaned dataset contained 97 cartoon-animal images, 97 plush-animal images, and 95 real-animal images (289 total).

% ---------------------------------------------------------------------------
\section*{Problem 2. Preparing an image}

\sidebyside{figures/prepare_original.jpg}{\fillin{the image as downloaded}}
           {figures/prepare_result.png}{\fillin{what prepare\_image made of it}}

\fillin{What was cut off, and an image where the crop cut off part of the
thing you care about.}

% ---------------------------------------------------------------------------
\section*{Problem 3. First training runs}

\begin{center}\footnotesize
\begin{tabular}{llllrrrr}
\toprule
tag & start & trained & epochs & lr & trainable parameters & train acc. & test acc. \\
\midrule
run & ImageNet & all layers & 10 & 0.0001 & 11,178,051 & 100.0\% & 93.0\% \\
scratch & random & all layers & 10 & 0.0001 & 11,178,051 & 84.9\% & 64.9\% \\
frozen & ImageNet & last layer & 10 & 0.001 & 1,539 & 99.6\% & 94.7\% \\
longer & ImageNet & all layers & 20 & 0.0003 & 11,178,051 & 97.4\% & 84.2\% \\
\bottomrule
\end{tabular}
\end{center}

\picturefile{figures/run_curves.png}

Training loss decreased in all four runs, but the test results did not improve steadily. In the longer run, test accuracy fluctuated and ended at 84.2%, below the 93.0% of the default run. More epochs and a larger learning rate therefore did not guarantee better generalization.

The default run reached 100.0% training accuracy and 93.0% test accuracy, a 7.0 percentage-point gap. The frozen run reached 99.6% training and 94.7% test accuracy, a 4.9-point gap. These gaps suggest some overfitting, while the very high training accuracy shows that the models learned the training examples much more easily than unseen test images.

\begin{center}\picturefile[0.5\linewidth]{figures/run_confusion.png}\end{center}

The confusion matrix should be inspected to identify the most frequent pair of mixed-up categories. I have not written a category-specific conclusion here without verifying the matrix values.

\fillin{How far apart the test accuracies are, starting from ImageNet and from
random numbers, and why you think that is.}

Training only the last layer scored 94.7%, 1.7 percentage points above the default all-layer run (93.0%). For this dataset, the pretrained features were already useful enough that updating only the classifier performed competitively. The longer all-layer run reached only 84.2%, illustrating that more training was not automatically better.

\subsection*{Output of the runs}
\outputfile{outputs/run_output.txt}
\outputfile{outputs/scratch_output.txt}
\outputfile{outputs/frozen_output.txt}
% One more line like these for the setting you chose, named by its tag.
\outputfile{outputs/longer_output.txt}

% ---------------------------------------------------------------------------
\section*{Problem 4. Cleaning}

\subsection*{My rule}
\fillin{Your rule, in one or two sentences, as you wrote it before you
started.}

\subsection*{Three kinds of junk}
The recorded removal reasons show three kinds of unsuitable images: (1) collages or illustrations instead of real-animal photos, (2) real animals or toys incorrectly labelled as cartoon animals, and (3) real-animal photos incorrectly labelled as plush toys. The individual example pictures still need to be inserted.

% For example: \picturefile[0.3\linewidth]{figures/junk_1.jpg}

\subsection*{What I removed}
\outputfile{outputs/clean_count.txt}

The checker reported one duplicate image within the same class. The duplicate warning should be resolved or explicitly documented before submission; the current check output does not establish that it was removed.

\fillin{Of the first 20 suspects in each class, how many you removed. Whether
the suspects included junk you had missed, and whether you found junk it did
not list.}

\subsection*{Before and after cleaning}
\textit{Comparison not run: the available outputs contain no before/after comparison. The initial training runs were mistakenly performed on the cleaned dataset, so a valid comparison requires retraining the baseline on the original dataset.}

\fillin{The test accuracy and the accuracy on your new images, before and
after. If a number went down, why you think it did.}

% ---------------------------------------------------------------------------
\section*{Problem 5. Images from a new source}

\fillin{Where your new images came from, and why you are sure none of them are
among your downloads.}

\textit{Overlap check not run because data/new_images/ does not exist.}

\fillin{The accuracy on your downloaded test set and on your new images, side
by side.}

\fillin{Five mistakes it was confident about, with the pictures. For each one,
what in the image you think pushed it the wrong way.}

\fillin{What is different about your new images: background, lighting, angle,
distance, what else is in the frame.}

% ---------------------------------------------------------------------------
\section*{Problem 6. The demo}

Address: \url{https://haey03.github.io/mas1004-assignment1/}

\fillin{What happens when you show it something that is none of your
categories.}

% ---------------------------------------------------------------------------
\section*{Problem 7. What the model looks at}

% Copy this subsection once for every guess.
\subsection*{Guess 1}

\fillin{The guess, and what in your images or your model's mistakes made you
think of it.}

No controlled image-editing experiment is recorded yet. Choose one image, change only the suspected feature (for example, background while keeping the animal unchanged), and compare the model's predictions.

My request to the agent, word for word:
\begin{quote}
ADD THE EXACT PROMPT YOU GAVE THE AGENT FOR THE IMAGE-EDITING EXPERIMENT
\end{quote}

% One \sidebyside for every image you changed.
\sidebyside{figures/guess1_original_1.jpg}{\fillin{original: the model's answer}}
           {figures/guess1_changed_1.jpg}{\fillin{changed: the model's answer}}

The hypothesis cannot be evaluated until the original and changed images have been tested with the model.

\subsection*{What I now think the model uses}
\fillin{What you now think the model uses to tell your categories apart, and
how sure you are.}

% ---------------------------------------------------------------------------
\section*{Appendix. Output of check.py}
\outputfile{outputs/check.txt}

\end{document}
